\chapter*{Conclusion et perspectives}
\addcontentsline{toc}{chapter}{Conclusion et perspectives}
\markboth{Conclusion et perspectives}{}

\section{Conclusion}

Le d\'eveloppement de la plateforme CyberGuard B\'enin a \'et\'e l'occasion de mettre
en pratique l'ensemble du cycle de vie d'un projet informatique~: clarification
du besoin, mod\'elisation UML, conception de la base de donn\'ees, puis
impl\'ementation technique. Bas\'ee sur une architecture combinant une API Laravel
et une interface React, la plateforme automatise un diagnostic de s\'ecurit\'e
jusque-l\`a r\'eserv\'e \`a des outils techniques et anglophones, en le rendant
accessible \`a un public non-technicien gr\^ace \`a un score, un niveau de risque et
des recommandations en fran\c{c}ais.

Au terme du stage, la plateforme livr\'ee couvre l'ensemble des cas d'utilisation
identifi\'es en phase de conception~: authentification (par mot de passe et via
Google), lancement et traitement asynchrone des analyses (sept contr\^oles non
intrusifs), consultation des rapports avec t\'el\'echargement PDF, tableau de
bord, notifications in-app de fin d'analyse, ainsi qu'un back-office
d'administration complet (gestion des comptes, mod\'eration des analyses,
param\'etrage du bar\`eme de scoring). Le principal point rest\'e en dehors du
p\'erim\`etre du stage est le d\'eploiement en production, la plateforme n'ayant
\'et\'e ex\'ecut\'ee qu'en environnement de d\'eveloppement local.

\section{Perspectives}

Parmi les perspectives d'\'evolution envisag\'ees~:
\begin{itemize}
  \item le d\'eploiement en production et l'ouverture de la plateforme au
    public, avec confirmation du HTTPS de bout en bout~;
  \item la v\'erification de l'adresse email \`a l'inscription (compte cr\'e\'e par
    formulaire classique)~;
  \item le suivi automatique p\'eriodique~: relancer r\'eguli\`erement l'analyse
    d'un site d\'ej\`a enregistr\'e et alerter par email en cas de r\'egression du
    score, plut\^ot que de d\'ependre uniquement d'une analyse manuelle~;
  \item un rapport public partageable (lien en lecture seule) et un badge de
    score int\'egrable sur le site audit\'e~;
  \item l'ajout de v\'erifications suppl\'ementaires (CSP plus fine, en-t\^etes
    additionnels, historique des CVE connues li\'ees aux technologies
    d\'etect\'ees)~;
  \item une couverture de tests automatis\'es \'etendue au frontend, en
    compl\'ement de la suite PHPUnit existante c\^ot\'e backend.
\end{itemize}

En d\'efinitive, CyberGuard B\'enin se veut une contribution concr\`ete \`a la
d\'emocratisation de la cybers\'ecurit\'e au B\'enin, en r\'eduisant la surface
d'attaque des sites web mal s\'ecuris\'es et en rendant la premi\`ere \'etape d'un
audit de s\'ecurit\'e accessible \`a tous.
